\chapter{Testes de Verificação e Validação}

\section{Características Gerais}

\textcolor{red}{Introduzir os principais pontos deste capítulo. Cada teste deve conter explicações completas que garantam sua repetibilidade:
\begin{itemize}
    \item Hipóteses levantadas
    \item Condições de contorno
    \item Resultados esperados
    \item Materiais e métodos
    \item Precisão e acurácia das medidas obtidas.
\end{itemize}
}

\section{Testes da estrutura}

\textcolor{red}{Apresentar com detalhes os testes feitos para conferir se os objetivos da estrutura foram atendidos, incluindo:}

\begin{itemize}
    \item \textcolor{red}{Imagem com todas as peças separadas da estrutura (carcaça) do veículo;}
    \item \textcolor{red}{Imagem que apresente a estrutura já montada (sem necessidade de partes internas, como eletrônica ou motores);}
    \item \textcolor{red}{Ilustração contemplando o carro montado virtualmente em alguma ferramenta de CAD, no intuito de verificar se as peças foram construídas e montadas em conformidade com o projeto.}
\end{itemize}

\section{Testes de \textit{hardware}}

\textcolor{red}{Apresentar com detalhes os testes feitos para conferir se os objetivos do \textit{hardware} foram atendidos, incluindo:}

\begin{enumerate}
    \item \textcolor{red}{Processamento (Arduino, ESP32 etc):}
    \begin{itemize}
        \item \textcolor{red}{Definir interface de desenvolvimento;}
        \item \textcolor{red}{Rodar exemplo mínimo.}
    \end{itemize}
    \item \textcolor{red}{Sensores (temperatura, pressão, posição, velocidade etc.):}
    \begin{itemize}
        \item \textcolor{red}{Realizar leitura básica de dados;}
        \item \textcolor{red}{Validar/calibrar leituras;}
        \item \textcolor{red}{Definir taxas de amostragem e quantidade de dados para transmissão e/ou armazenamento;}
        \item \textcolor{red}{Validar taxas definidas e quantidade de dados;}
        \item \textcolor{red}{Realizar processamento básico de dados (\textit{thresholding}, filtro média-móvel etc.).}
    \end{itemize}
    \item \textcolor{red}{Armazenamento (memória interna, cartão SD, pendrive etc.):}
    \begin{itemize}
        \item \textcolor{red}{Armazenar dados básicos;}
        \item \textcolor{red}{Armazenar e validar dados de sensores.}
    \end{itemize}
    \item \textcolor{red}{Atuadores (Motores DC, relés, LEDs etc.):}
    \begin{itemize}
        \item \textcolor{red}{Realizar acionamentos básicos;}
        \item \textcolor{red}{Realizar acionamentos com base em dados de sensores;}
        \item \textcolor{red}{Realizar processamento básico de dados (PWM etc.).}
    \end{itemize}
    \item \textcolor{red}{Comunicações (UART, I2C, SPI, USB, WiFi, Bluetooth etc.):}
    \begin{itemize}
        \item \textcolor{red}{Transmissão básica de dados;}
        \item \textcolor{red}{Transmissão de dados de sensores dentro da taxa de amostragem.}
    \end{itemize}
\end{enumerate}

\section{Testes de consumo energético}

\textcolor{red}{Apresentar com detalhes os testes feitos para conferir se as demandas energéticas do projeto foram atendidas.}

\section{Testes de \textit{software}}

O \textit{backlog} do produto definido no primeiro ponto de controle (PC1) é composto por seis histórias
de usuário (HU01--HU06). Nesta entrega foram efetivamente implementadas e validadas
por testes automatizados as quatro histórias do módulo de \textit{software}
(web/\textit{backend}): HU03, HU04, HU05 e HU06, o que corresponde a $66{,}7\%$ do
\textit{backlog} funcional (acima do mínimo de $50\%$). As histórias HU01
(resolução autônoma) e HU02 (mapeamento incremental) pertencem ao módulo embarcado
e dependem da integração final com a placa ESP32, estando planejadas para a etapa
seguinte; por isso não compõem o conjunto validado nesta seção.

\begin{table}[H]
    \centering
    \caption{Associação entre histórias de usuário implementadas, protótipos e
    funcionalidades acessadas na interface gráfica.}
    \label{tab:hu_prototipo_gui}
    \footnotesize
    \begin{tabular}{|c|p{3.6cm}|c|p{3.6cm}|p{3.4cm}|}
        \hline
        \textbf{HU} & \textbf{História de usuário} & \textbf{Protótipo} &
        \textbf{Funcionalidade acessada na GUI} &
        \textbf{Suíte de testes que valida} \\ \hline

        HU03 & Telemetria em tempo real: trajeto, tempo, velocidade média e
        bateria durante a corrida. & Figura~\ref{fig:telemetria_fig} &
        Tela \textit{Telemetria ao Vivo}: grade do labirinto,
        métricas e status de conexão atualizados via WebSocket. &
        Unitário (\textit{Telemetry}, \textit{useTelemetry}); E2E
        \texttt{telemetria.spec.js}; \textit{backend} POST \texttt{/telemetria}
        e WS \texttt{/ws/telemetria}. \\ \hline

        HU04 & Persistência e consulta de histórico de corridas com filtro por
        tamanho. & Figura~\ref{fig:historicocorridas_fig} &
        Tela \textit{Histórico}: tabela de corridas
        e filtro por dimensão do labirinto. &
        Unitário (\textit{History}, \textit{useHistory}); E2E
        \texttt{history.spec.js}; \textit{backend} GET \texttt{/corridas}. \\ \hline

        HU05 & Monitoramento do nível de bateria no painel web. &
        Figura~\ref{fig:monitoramentobateria_fig} &
        Indicador de bateria na tela \textit{Telemetria ao Vivo}, alimentado pelo 
        campo \texttt{nivel\_bateria} da telemetria. &
        Coberta pelos testes de telemetria (unitário e E2E); validada com dados
        simulados. \\ \hline

        HU06 & Indicador de desafio cumprido (Sim/Não) ao atingir o centro do
        labirinto. & Figura~\ref{fig:desafiocumprido_fig} &
        Campo \textit{Desafio Cumprido} na tela \textit{Telemetria ao Vivo}, 
        calculado ao atingir a região central. &
        Unitário (\textit{Telemetry}); E2E \texttt{telemetria.spec.js}
        (casos ``Sim'' e ``Não''). \\ \hline
    \end{tabular}
\end{table}

\textcolor{red}{Apresentar com detalhes os testes feitos para conferir se os objetivos do \textit{software} foram atendidos.
Considerando a implementação de $50\%$ do \textit{backlog} do produto (funcional):}

\begin{itemize}
    \item \textcolor{red}{Listar as histórias de usuários implementadas.}
    \item \textcolor{red}{Listar a associação entre as histórias de usuários e os protótipos, deixando claro qual é a funcionalidade que será acessada pelo usuário/professor(a) por meio da interface gráfica.}
    \item \textcolor{red}{Listar os suítes de testes automatizados em nível:}
    \begin{itemize}
        \item \textcolor{red}{Unitário e integração, com cobertura de código-fonte de pelo menos $70\%$;}
        \item \textcolor{red}{Sistema, do tipo funcional (end-to-end E2E).}
    \end{itemize}
\end{itemize}

\textcolor{red}{\textbf{Observações:}}
\begin{enumerate}
    \item \textcolor{red}{Em relação ao \textit{backlog} do produto, será considerado como referência aquele descrito no primeiro ponto de controle, contemplando as correções e \textit{feedbacks} da-os professora(e)s. O ideal seria a real atualização entre o planejado X realizado.}
    \item \textcolor{red}{Os testes mencionados nesta subseção dizem respeito aos testes de \textit{software}. O Roteiro de testes funcionais, entregue no primeiro ponto de controle deve ser corrigido a partir do apontamento da-os professora(e)s. Além disso, ele deve guiar a execução dos testes E2E automatizados. O conjunto de casos de testes funcionais documentados devem estar coerentes com a automação via interface gráfica de entrada de dados.}
\end{enumerate}

\textcolor{red}{Exemplos de ferramentas de teste funcionais E2E:}
\begin{itemize}
    \item \textcolor{red}{Selenium} 
    \item \textcolor{red}{Cypress}
    \item \textcolor{red}{Playwright}
\end{itemize}

\textcolor{red}{Exemplos de ferramentas de testes em nível unitário e integração:}
\begin{itemize}
    \item \textcolor{red}{Jest}
    \item \textcolor{red}{unittest}
    \item \textcolor{red}{pytest}
    \item \textcolor{red}{tox}
    \item \textcolor{red}{gTest}
    \item \textcolor{red}{Catch2}
\end{itemize}

\section{Testes de integração}

\textcolor{red}{Apresentar com detalhes os testes feitos para conferir se os objetivos do projeto como um todo foram atendidos.}

%\textcolor{red}{Com relação ao \textit{software}, será necessário apresentar pacotes de componentes de \textit{software}, suas funções e características, e explicar as decisões de projeto.}

